\documentclass[11pt,a4paper]{article}

% --- Pacchetti ---
\usepackage[utf8]{inputenc}
\usepackage[T1]{fontenc}
\usepackage{lmodern}
\usepackage[italian]{babel}
\usepackage[margin=1in]{geometry}
\usepackage{xcolor}
\usepackage{listings}
\usepackage{hyperref}
\usepackage{enumitem}

% --- Colori ---
\definecolor{primary}{RGB}{38, 66, 139}
\definecolor{codebg}{RGB}{245, 245, 245}

\hypersetup{colorlinks=true, linkcolor=primary, urlcolor=primary}

% --- Configurazione Codice ---
\lstdefinelanguage{javascript}{
    keywords={async, await, break, case, catch, class, const, continue, debugger, default, delete, do, else, export, extends, false, finally, for, function, if, import, in, instanceof, new, null, return, super, switch, this, throw, true, try, typeof, var, void, while, with, yield, let},
    morekeywords={console, log, Error, fetch, JSON, parse, stringify},
    sensitive=true,
    morecomment=[l]{//},
    morecomment=[s]{/*}{*/},
    morestring=[b]',
    morestring=[b]",
    morestring=[b]`,
    keywordstyle=\color{primary}\bfseries,
    commentstyle=\color{green!50!black}\itshape,
    stringstyle=\color{red!70!black},
    basicstyle=\ttfamily\small
}

\lstset{
    backgroundcolor=\color{codebg},
    basicstyle=\ttfamily\small,
    breaklines=true,
    frame=single,
    keywordstyle=\color{blue},
    commentstyle=\color{green!50!black},
    stringstyle=\color{red!70!black},
    language=javascript
}

\title{\textbf{\color{primary} Esercitazione: Sviluppo di un'API REST Full-Stack}}
\author{Programmazione e Tecnologie Web - Lezione 4}
\date{A.A. 2025/2026}

\begin{document}

\maketitle

\section{Obiettivo}
L'obiettivo di questa esercitazione è trasformare una Todo App statica in un sistema \textbf{Full-Stack} completo. Impareremo a gestire il ciclo di vita delle risorse su un server, implementare la persistenza dati e garantire la sicurezza tramite \textbf{Basic Authentication}.

\section{Prerequisiti}
\begin{itemize}
    \item Ambiente di runtime: \textbf{Node.js} (v18+) o \textbf{Python} (3.10+).
    \item Estensione VS Code: \textbf{Thunder Client} o \textbf{REST Client}.
    \item Codice sorgente nella cartella \texttt{src/lezione4/1.basic-todo/}.
\end{itemize}

\section{Esercizio 0: Analisi e Setup}
Esplorate il codice fornito nella cartella \texttt{1.basic-todo}. Noterete che il backend (\texttt{app.js} o \texttt{app.py}) carica i dati dal file \texttt{tasks.json} solo all'avvio.

\textbf{Guida per lo studente}: 
\begin{enumerate}
    \item Aprite il terminale nella cartella \texttt{server/}.
    \item Avviate il server:
    \begin{itemize}
        \item Node.js: \texttt{npm install && npx nodemon app.js}
        \item Python: \texttt{python app.py}
    \end{itemize}
    \item Aprite l'app nel browser (es. tramite Live Server o via server backend su \texttt{http://localhost:3000}).
\end{enumerate}

\section{Esercizio 1: La Persistenza "Volatile"}
Attualmente, se aggiungete un task, questo viene salvato solo nella memoria RAM del server.
\begin{itemize}
    \item \textbf{Test}: Aggiungete un task dall'interfaccia. Ricaricate la pagina: il task c'è ancora (perché è in memoria).
    \item \textbf{Problema}: Riavviate il server (es. premete \texttt{Ctrl+C} e riavviate). Ricaricate la pagina: il task è sparito!
\end{itemize}

\textbf{Task Live}: Identificate nel codice del server il punto in cui i dati vengono letti dal file e spiegate perché la funzione \texttt{add\_task} (o la rotta \texttt{POST}) non garantisce la persistenza a lungo termine.

\section{Esercizio 2: Implementazione della Persistenza Reale}
Vogliamo dare all'utente il controllo su quando salvare i dati sul disco.

\subsection{Task Backend: Rotte /save e /reset}
Aggiungete due nuovi endpoint al vostro server:
\begin{enumerate}
    \item \texttt{POST /save}: Deve scrivere l'array corrente di task nel file \texttt{tasks.json}.
    \item \texttt{POST /reset}: Deve svuotare l'array in memoria e sovrascrivere il file \texttt{tasks.json} con un array vuoto \texttt{[]}.
\end{enumerate}

\begin{lstlisting}[language=javascript, caption=Esempio Express]
app.post('/save', (req, res) => {
    fs.writeFileSync(FILE_PATH, JSON.stringify(tasks, null, 4));
    res.json({ status: "saved" });
});
\end{lstlisting}

\section{Esercizio 3: UI Interattiva}
Modificate il file \texttt{index.html} per includere i comandi di gestione persistenza.

\textbf{Task}:
\begin{enumerate}
    \item Aggiungete due pulsanti sotto la lista dei task: "Salva su Disco" e "Resetta Tutto".
    \item Create due funzioni JavaScript \texttt{saveTasks()} e \texttt{resetTasks()} che utilizzano \texttt{fetch} per chiamare i nuovi endpoint.
    \item Aggiornate la UI dopo il reset.
\end{enumerate}

\begin{lstlisting}[language=javascript]
async function saveTasks() {
    await fetchAPI('/save', 'POST');
    alert('Dati salvati con successo!');
}
\end{lstlisting}

\section{Esercizio 4: UX Feedback e Latenza Simulata}
In una rete reale, le richieste API non sono istantanee. È fondamentale fornire un feedback visivo all'utente per evitare che l'interfaccia sembri "bloccata".

\subsection{Task Backend: Simulare la lentezza}
Per testare il caricamento, forziamo il server a rispondere dopo un ritardo (es. 2 secondi).

\begin{lstlisting}[language=javascript, caption=Ritardo in Node.js]
// Inserire all'inizio della rotta GET /tasks
await new Promise(resolve => setTimeout(resolve, 2000));
\end{lstlisting}

\begin{lstlisting}[language=Python, caption=Ritardo in Python]
import time
# Inserire nella rotta get_tasks()
time.sleep(2)
\end{lstlisting}

\subsection{Task Frontend: Lo Spinner di caricamento}
\begin{enumerate}
    \item Aggiungete nel file HTML un elemento \texttt{<span>} o un div (es. con ID \texttt{loader}) che contenga il testo "Caricamento in corso..." e uno spinner CSS.
    \item Gestite la visibilità del loader nella funzione \texttt{render()}:
    \begin{itemize}
        \item Mostrate il loader appena inizia la funzione.
        \item Nascondetelo subito dopo aver ricevuto i dati dal server.
    \end{itemize}
\end{enumerate}

\begin{lstlisting}[language=javascript, caption=Gestione Stato nel Frontend]
async function render() {
    loader.classList.remove('hidden'); // Mostra
    const tasks = await fetchAPI('/tasks');
    loader.classList.add('hidden');    // Nascondi
    // ... resto della logica di rendering
}
\end{lstlisting}
\begin{enumerate}
    \item \textbf{Gestione Errori}: Modificate la \texttt{fetchAPI} per mostrare un messaggio di errore all'utente se il server restituisce uno stato \texttt{500} o \texttt{401}.
    \item \textbf{Audit Log}: Ogni volta che un task viene eliminato, stampate nel terminale del server il timestamp e il testo del task rimosso.
    \item \textbf{Task Live (Avanzato)}: Implementate una rotta \texttt{GET /stats} che restituisca il numero totale di task e la lunghezza media del testo dei task. Visualizzate questi dati in un piccolo box nell'interfaccia.
\end{enumerate}

\end{document}
